\section{Paper Checklist}\label{apdx:paper checklist}

Here we explicitly comment on all areas of the NeurIPS paper checklist.

\paragraph{Claims}
The key claims in this paper are the creation of random features for the Tanimoto kernel and the creation of a new continuous kernel.
The claims are mainly supported theoretically with proofs,
but we also test our random features experimentally and show that they are usable with real-world data.

\paragraph{Code of Ethics + broader impacts}
Read and acknowledged.
Our work is fairly abstract and theoretical and we do not work with human-related data,
so we do not foresee significant direct ethical impacts
(positive or negative).

\paragraph{Limitations}
Our work is primarily theoretical, so the main strength (random features with low approximation error) is also the main weakness (the approximation error could plausibly be lower).
Random features are well-studied as a general approach so the strengths and weaknesses are generally well-understood \citep{liu2021random}.
For this reason we did not include an explicit limitations section.
However, we did mention some other limitations in section~\ref{sec:conclusion}.

\paragraph{Theory}
Our key theorems have minimal assumptions which are clearly stated.
All theorems are proved in Appendix~\ref{proofs in appendix}.

\paragraph{Experiments}
Code to reproduce the experiments is available at:\newline
\codeURL{}.

\paragraph{Training details}
The experiments in this paper have minimal details and the important details are specified in Appendix~\ref{appendix experimental details}.
All unspecified details should be easy to find in the code (\codeURL{}).

\paragraph{Error bars} Our tables and figures include error bars.

\paragraph{Compute}
The compute costs of the experiments in this paper were quite modest
and were run on a single machine with no GPU usage.
The approximate computational times were:
\begin{itemize}
    \item Section~\ref{sec:expt-rf-on-real-data}: a script to produce all the plots took approximately 24\,h to run.
    The long runtime was mostly due to computing errors for many settings of $R$ and $p$.
    \item Section~\ref{sec:regression-expt}:
        for $T_{MM}$, SVGP and each RFGP inference took approximately 1\,h (total 3\,h per target).
        For $T_{DP}$, SVGP took 0.5\,h while RFGP took 0.3\,h (total 1.5\,h per target).
        Multiplying this by 5 targets and 5 trials gives a total compute time of $\approx$112\,h.
    \item Section~\ref{sec:expt-bayesopt}:
        this experiment was quite fast: total runtime was perhaps 5\,h.
\end{itemize}

\paragraph{Reproducibility}
Our contribution is reproducible via the statements of our random features
and our code.
We aim for a high standard of reproducibility:
the exact commands to reproduce the results are stated explicitly in our code's \texttt{README}
file,
and the raw data behind all plots in this paper is included alongside the code.

\paragraph{Safeguards}
We believe there is nothing high-risk which we need to safeguard.

\paragraph{Licenses}
We cite the assets which we use in the paper.

\paragraph{Assets}
We are not releasing assets.

\paragraph{Humans Subjects / IRB Approvals}
Not applicable to our paper.
